\section{Letter to Palamidessi}

Rome, January 20, 1972

Dear Dr. Palamidessi,

I got your letter, for which I thank you.

I already told you that a treatise on initiatic astrology would be an important work and would fill a gap; and it would also be interesting to treat, in accordance with its aim, a typology of inner races with the utilization of the astrological element. In my book, there is only a sketch of it that needs to be adequately developed. A publisher in Germany had encouraged me to dedicate myself to this research, but I had to pass it up. Moreover, it would, above all, be about an \emph{a priori} morphology, not on an empirical base. Because of that, whether pure types currently exist, conformed completely to that fundamental structure, has no importance, and even lacks practical value. So, for example, the defining of the absolute woman and the absolute man, who are almost non-existent empirically, is of great utility for the analysis of those women and men who are such only approximately, and with admixtures (Weininger).

The temperamental knowledge of oneself in the deepest states would be an important possibility offered by such investigations. I don't see, on the other hand, possible connections to that ``awakening of the European race'', which you noted. An earlier work, although without utilizing any astrological contributions, could be my friend L. W. Clauss' Rassenseelekunde racial psychology (the study of the soul or the souls of race); I believe that I mentioned him in my book. If you want, you can include me in those figure sketches that you intend to treat, provided that you exclude every exaggeration, as the English would say. One of my accounts of the ``absolute individual'', for example, would only be humorous.

Unfortunately, I forgot to give you directly, on the occasion of your visit, the book L'Arco e la Clava. We will see how it can be done: either on the next meeting or if the publisher provides the service of sending right to a private home, the postal services particularly for the stamps have become infamous. However, a copy is reserved for you.

Another thing. One of my (female) friends, belonging to my faithful guard-house, having known of your visit and knowing your name, makes my life impossible because she would like to have from you an astrological chart of her character and her ``destiny''. If, by chance, you have a little time to waste for a little more than summary investigation, the dates are: born February 25, 1945 at 15:20 in Venice.

Cordially yours,

J. Evola


\flrightit{Posted on 2012-04-30 by Cologero}

\begin{center}* * *\end{center}

\begin{footnotesize}
\begin{sffamily}
\texttt{Angolmois on 2012-04-30 at 01:48 said:}

Thanks for this one, it made my day, haha!

Seriously, when reading Evola and Guénon and others of their ilk only from their theoretical writings, one gets easily the impression that they are some kind of mythological heroes who perhaps never even lived in the first place. When one sees Evola discussing postal costs of stamps and talking in an embarrasing way about astrological charts etc., it kind of brings him as a human being more down to earth without making his writings any less lofty. I can only imagine him trading tapes of rare traditionalist BM bands with his UG friends. :-)

\hfill

\texttt{Matt on 2012-04-30 at 07:10 said:}

Yes, I thought the same thing when I watched the video of Evola being interviewed on the history of Italy's Dada movement. Evola the author comes across as being aloof and distant, but Evola in-person comes across as quite friendly and close.

\hfill

\texttt{Cologero on 2012-04-30 at 07:28 said:}

Yes, Angolmois, every man puts on his pants, one leg at a time. The caretaker story is curious, too, that he would make the effort to get a private astrological reading for her.

But it also shows his consistency over the course of his life, unlike many men I know who go through a major conversion every few years. He still refers, in a positive way, to Weininger's \textit{Sex and Character}\footnote{\url{http://www.gornahoor.net/library/SexAndCharacter.pdf}}, which was one of his earliest influences. I know Evola ``scholars'' who prefer to ignore Weininger.

We get an insight, too, about Evola's method and way of thinking. The Absolute Man or Woman is an \emph{a priori} metaphysical category; actually existing men and women are to be compared to that standard. This is the opposite approach of empirical science, which would first do a study of men and women as they are; then, they would formulate a theory of what is a man or a woman. It seems to me that post-modern science has brought much more confusion than clarity to that question.

The methods are clearly incompatible and lead to much different results. Note, too, that those on the ``right wing'' today try to employ biological arguments to demonstrate intrinsic differences between men and women. This approach has proved less than effective as it falls far short of an intellectual conversion to a metaphysical and transcendental world view.

There are other questions, if you read between the lines. Why would Evola be embarrassed by his writings on the absolute individual, one of very earliest conceptions? There is also the reference in passing to the relationship of astrology to race; was Evola thinking that had something to do with the problem of births?

Perhaps Logres can shed some light on these questions:

\begin{enumerate}


\item Did P. ever publish his sketch of Evola?

\item Did P. ever publish anything on ``awakening the European race''?
\end{enumerate}


\hfill

\texttt{Angolmois on 2012-05-01 at 05:41 said:}

Good points, both of you. The concept of the absolute man brought to my mind Guénon's ``nickname'' which was `an intellectual function'. What is man and how we should jugde him according to his accidental traits is certainly a lost call, yet in differentiated men it does also shed some light upon the personality behind the reflection and the mask / per-sona.

Evola liked talking about his `personal equation' although he scorned his own personality to the borders of sanity. Quite telling again, isn't it? We all have our shadows to deal with.

\hfill

\texttt{logres on 2012-05-01 at 08:14 said:}

This is probably the book of interest: \textit{Astrologia mondiale: il destino dei popoli rivelato dal corso degli astri}, Turin: T. Palamidessi, 1941 (2nd ed. Archeosofica 1985).

(The destiny of people revealed by the course of the stars, I think). I'll have to look for more. Many thanks to Cologero for translating this letter, which has piqued an interest in some of his posts on cosmology and the zodiac (although there are fewer of them). ``We are underlings to the stars'', except for our will, if it awakens.

\hfill

\texttt{Charlotte on 2012-05-01 at 08:27 said:}

the astrological tangent puts me in mind of the words of a truly catholic author along the lines of Justin or Origen, who wrote something along the lines of (quoted in MoTT) -- `When we were children, under compulsion from the stars'.

This has always struck me as one of the most poignant phrases I've read in any book anywhere -- the innocence of the child who is both protected and controlled by these cosmic forces, while the adult man or woman must master those same forces in order to attain the truth contained in another wonderful phrase from MoTT -- the stars incline, they do not compel.

That probably tells us all we need to know about astrological sciences, which certainly doesn't stop me gazing at the planets and stars as often as possible, especially Venus, which is monumentally huge right now -- the biggest it will EVER appear to us in fact, so far as I know (last night) :-)

\hfill

\texttt{Angolmois on 2012-05-03 at 01:35 said:}

Charlotte: I think that the two influences of Venus and Jupiter, and later on Saturn, can be traced to the beginnings of the new age in astrological terms. The conjunction of Venus and Jupiter marks the birth of the Christ, and they have both been active during this year. There is also the issue of the comet that stroke the sky into the direction of the constellation Draco which marked the birth of the year of the Dragon. We are currently witnessing the inauguration of the reign of Saturn and the age of Aquarius; this is the very thing in which the new agers are certainly right, only their conclusions and the applications seem to be often quite fantastic and utopian. The major star Betelgeuse has also activated, which can be linked to the emergence of `the true will' in highly developed humanity.

Personally I have found the Bailey book on esoteric astrology quite helpful, although I can't be completely in agreement with Bailey's conclusions or theories.

As regards the `enslavement under the stars', this is mostly a problem for mass minded humanity. For every disciple and initiate the birth chart and astrological researches must be made according to their point in the path of return together with other issues to be examined, to be able to tell one's exact point.

In whole I think that the (esoteric) teachings and interpretations of Jesus Christ are the most helpful in ``surmounting the archons'', especially the teachings found from the apocryphic writings and Pistis Sophia, which might just be the most neglected ``christic'' work of all times.

\hfill

\texttt{Charlotte on 2012-05-03 at 06:30 said:}

yes Angolmois, I spent a fascinating month in the highlands of Guatemala, from a wonderful natural observatory, watching the conjunctions, comets and other incredible astronomical phenomena you mention. Betelgeuse does indeed appear very red nowadays (all the more obvious in a clear sky) and it struck me for the first time that it might indeed go (or already be going) supernova as expected. Imagine Orion without his left arm! I totally agree that the cosmos has been behaving in an unusual way -- or a way not seen in my liftime before -- since at least last September, when I saw Jupiter looking so enormous it stunned me. Now Venus is doing the same, and as you pointed out, they have already conjuncted. I have mentioned to people a few times that early in October I saw literally the most amazing night sky I've ever witnessed in my life, it had both myself and my neighbour out on the balcony with our jaws on the floor. It only lasted for about 10 minutes, but for that time every single star and planet seemed about 10 times closer, you could see every colour -- red, blue, green, orange -- and they weren't just twinkling in the usual fashion, Jupiter for instance was slowly pulsating as if it was actually turning on and off. Anyone who denies there are signs in the skies these days either hasn't been looking up or they live in a zone where there's too much ambient light and air pollution to see properly. What it all means, as you also rightly pointed out, is a matter of interpretation. Clearly this IS the start of the age of Aquarius, that's just a fact connected with our cosmos... I am also with you 100\% on the Pistis Sophia and often recommend it to people. I broke my foot last January (over a year ago) , tripping over in towering heels on a slippy floor, and it gave me perfect occasion to contemplate the fall of man and womankind with the most tragic of attitudes, I can tell you! Joking aside, this was a very dark and difficult period of my life that caused me to examine myself thoroughly and face up to things that hitherto I'd been unwilling or unable to deal with, but reading the Pistis Sophia for the first time was definitely a silver lining in this cloud... .it is an incredible work, it is of so much help to anyone contemplating the problem of the Fall, a truly magical text. I posted this extract on my blog at the time -- looking back now I can see that I began reading it about 3 days after my tumble, something must have been guiding me for me to have found it soo soon after!

\url{http://alchemical-weddings.com/alchemical-weddings/the-ineffable}


Back to the sky for a moment, we have a red supermoon in Scorpio coming on Saturdady, I'd say that's a biggie. It's the Beltane full moon and also around the time of Buddha's enlightenment... .got to be a good time to meditate by any estimation?! Cx

\hfill

\texttt{Charlotte on 2012-05-03 at 06:35 said:}

yes it is curious that Pistis Sohpia has been so neglected... ..just like Mary Magdalene in exoteric religion

\hfill

\texttt{Angolmois on 2012-05-03 at 13:27 said:}

The neglection of Pistis Sophia among `orthodox' believers and exoteric christianity must have something to do with the `egyptian' and other influences in it. Amenti etc. Early Christians around the years 100-300 did quite a good job of kicking off almost all universal esoterism and twisting the words of Jesus, if I'm right. It all started from `Paul the theologian' and `Peter the denier'. Origen and few others are quite good, but even they saw gnosticism as their enemy becase they thought it leads men towards over-blown intellectual pride. On the other hand they were right, but on the other hand they were quite wrong.

AS regarding Mary Magdalene, I think this neglection has to do with too much Solarism understood in a unilateral way among the early christians and the desert fathers. The Eternal Feminine has certainly been neglected in the exoteric monotheism, yet its resurgence can be seen in the `wholly pagan' worship of the black virgin.

\url{http://www.sophiaperennis.com/books/christianity/the-black-virgin/}


\hfill

\texttt{Charlotte on 2012-05-03 at 15:14 said:}

Yes... .

The book sounds very interesting thank you for the link. In fact I was thinking of the Black Virgin in the bath earlier today, wondering if I should mention this. I had tended to see her as more equivalent to Mary Magdalene, but who knows?

The negation of Mary Magdalene -- or marginalisation, perhaps -- is definitely one of the main reasons I have felt unable to fully sign up to the catholic church. Peter and Paul have an awful lot of truth between them, but not a monopoly, although catholicism does more for the divine feminine than most other exoteric Christian denominations.

Funny you mentioned Amenti... .you know I keep on discovering things that are meant to be taboo! It reminds me of a poem or something else I have read, would that be the Emerald Tablet perhaps. Could be. Halls of Amenti. We've been there haven't we, isn't that where one of the schoolrooms is?!

\hfill

\end{sffamily}
\end{footnotesize}

